% ============================================================
%  Chapter 7 — Autonomous equations
% ============================================================
\chapter{Autonomous Equations and the Phase Line}\label{ch:ch07}

\begin{diagnostic}
  \item Make a sign chart for $f(x) = x(1-x)(x-2)$.
  \item State the mean value theorem, and use it to show: if $g$ is
        differentiable on $[a,\infty)$ and $g'(t)\ge c>0$ there, then
        $g(t)\to\infty$.
  \item If $f$ is differentiable at $a$ and $f(a) = 0$, $f'(a) = -3$, what is
        the sign of $f(x)$ for $x$ slightly larger than $a$? Justify from the
        definition of the derivative.
\end{diagnostic}

% ------------------------------------------------------------
\section{Motivation: how much fish can we take?}\label{sec:ch07-motivation}

A fish stock grows logistically and is harvested at a constant rate $h$
(tonnes per year):
\begin{equation}\label{eq:ch07-harvest}
  \der{x}{t} = rx\Bigl(1 - \frac xK\Bigr) - h .
\end{equation}
Managers want the largest $h$ that can be sustained forever. We could try to
solve \eqref{eq:ch07-harvest} (it is separable), but the formula depends on
$h$ in three different ways depending on the sign of a discriminant, and it
hides the one fact that matters: \emph{above a critical harvest, the
population collapses in finite time, whatever its starting size}. This
chapter develops a qualitative method that reads such facts directly from the
graph of the right-hand side.

\notation{Newton: in this chapter $t$ is time and we write $\dot x = \der{x}{t}$,
as is standard for dynamical systems. When we separate variables we switch back
to Leibniz and say so.}

% ------------------------------------------------------------
\section{Definitions}\label{sec:ch07-defs}

\begin{definition}[Autonomous equation; equilibrium]\label{def:ch07-autonomous}
An \emph{autonomous} first-order equation is $\dot x = f(x)$. A point $x^*$ with
$f(x^*) = 0$ is an \emph{equilibrium} (fixed point); the constant function
$x(t)\equiv x^*$ is a solution.
\end{definition}

Throughout this chapter (except where said otherwise) $f$ is $C^1$ on $\R$, so
every zero of $f$ satisfies the Lipschitz bound of
\cref{prop:ch03-lipschitz-zero} (a $C^1$ function is Lipschitz near each point,
by the mean value theorem).

\begin{proposition}[Time-translation invariance]\label{prop:ch07-translation}
If $x(t)$ solves $\dot x = f(x)$ on $(\alpha,\beta)$, then for every $c$,
$t\mapsto x(t - c)$ solves it on $(\alpha + c,\beta + c)$.
\end{proposition}
\begin{proof}
$\der{}{t}x(t-c) = \dot x(t-c) = f(x(t-c))$ by the chain rule. Note this uses
that $f$ does not depend on $t$.
\end{proof}

\begin{definition}[Stability]\label{def:ch07-stability}
An equilibrium $x^*$ is
\begin{itemize}
  \item \emph{stable} if for every $\varepsilon>0$ there is $\delta>0$ such that
        every solution with $|x(0) - x^*|<\delta$ exists for all $t\ge0$ and
        satisfies $|x(t) - x^*|<\varepsilon$ for all $t\ge 0$;
  \item \emph{asymptotically stable} if it is stable and, in addition, there is
        $\delta_0>0$ such that $|x(0) - x^*|<\delta_0$ implies $x(t)\to x^*$ as
        $t\to\infty$;
  \item \emph{unstable} if it is not stable.
\end{itemize}
\end{definition}

The definition of stability is a statement about \emph{all} nearby solutions
for \emph{all} future time; Chapter~20 uses the same definition in higher
dimensions.

% ------------------------------------------------------------
\section{The phase line, with proofs}\label{sec:ch07-theory}

\begin{theorem}[Structure of solutions on the line]\label{thm:ch07-structure}
Let $f$ be $C^1$ and let $x(t)$ be a solution of $\dot x = f(x)$ with
$f(x(t_0))\neq0$. Then:
\begin{enumerate}[label=(\alph*)]
  \item $x(t)$ never equals an equilibrium, and $x(t)$ is strictly monotone (increasing if $f(x(t_0))>0$, decreasing otherwise);
  \item if the solution moves toward a \emph{finite} endpoint $e$ of $J_0$ (the
        largest open interval containing $x(t_0)$ on which $f\neq0$), then it
        exists for all $t\ge t_0$ and $x(t)\to e$ as $t\to\infty$. In particular,
        if $x(t_0)$ lies between consecutive equilibria $a<b$, the solution
        exists for all $t\in\R$ and tends to one of $a,b$ as $t\to\infty$ and to
        the other as $t\to-\infty$;
  \item if $x(t)$ exists and is bounded on $[t_0,\infty)$, then $x(t)$ tends to an
        equilibrium as $t\to\infty$.
\end{enumerate}
\end{theorem}

\begin{proof}
Let $J_0$ be the largest open interval containing $x(t_0)$ on which $f\neq0$.

(a) Its finite endpoints are zeros of the $C^1$ function $f$, hence Lipschitz
zeros, so by \cref{cor:ch03-trapped} (with $g\equiv1$) the solution stays in
$J_0$. On $J_0$, $f$ has constant sign (intermediate value theorem), so
$\dot x = f(x(t))$ has constant sign, and $x$ is strictly monotone.

(b) Say $f>0$ on $J_0$, so the motion is to the right and $e$ is the right
endpoint (the other case is symmetric). Let $y_0 = x(t_0)$ and
$H(y) = \int_{y_0}^{y}\frac{\mathrm{d}s}{f(s)}$ for $y\in[y_0,e)$. Then $H$ is
increasing, and $H(y)\to+\infty$ as $y\to e^-$: since $f(e) = 0$ and $f$ is
Lipschitz near $e$, $0<f(s)\le L(e - s)$, so
$H(y)\ge\frac1L\ln\frac{e - y_1}{e - y}$ minus a constant, exactly as in the
proof of \cref{prop:ch03-lipschitz-zero}. Hence $H$ maps $[y_0,e)$ onto
$[0,\infty)$. In \cref{thm:ch03-separation}(c) we have $G(t) = t - t_0$, which
lies in $H(J_0)\supseteq[0,\infty)$ for all $t\ge t_0$; so the solution exists
for all $t\ge t_0$ and equals $H^{-1}(t - t_0)\to e$. If both endpoints are
finite equilibria, apply the same argument to the time-reversed equation
$\der{}{s}x(-s) = -f(x(-s))$, which moves toward the other endpoint.

(c) A bounded monotone function has a limit $L$ as $t\to\infty$. Suppose
$f(L)\neq0$, say $f(L)>0$. By continuity, $f(x(t))\ge \frac12 f(L)$ for all
$t\ge T$ (some $T$). By the mean value theorem,
$x(t) - x(T)\ge\frac12f(L)(t - T)\to\infty$, contradicting boundedness. So
$f(L) = 0$.
\end{proof}

\begin{corollary}[No oscillations in one dimension]\label{cor:ch07-no-periodic}
For $C^1$ $f$, $\dot x = f(x)$ has no nonconstant periodic solutions.
\end{corollary}
\begin{proof}
A nonconstant solution is strictly monotone by \cref{thm:ch07-structure}(a), so
it never repeats a value.
\end{proof}

\Cref{thm:ch07-structure} justifies the \emph{phase line}: draw the $x$-axis,
mark the zeros of $f$, and put an arrow pointing right where $f>0$ and left
where $f<0$. Every solution moves monotonically along an arrow and approaches
the equilibrium it points to. Whether a solution can reach $\pm\infty$ in
finite time is \emph{not} visible on the phase line
(\cref{ex:ch07-blowup}).

\begin{theorem}[Linearization test]\label{thm:ch07-linearization}
Let $f$ be $C^1$ and $f(x^*) = 0$.
\begin{enumerate}[label=(\alph*)]
  \item If $f'(x^*)<0$, $x^*$ is asymptotically stable.
  \item If $f'(x^*)>0$, $x^*$ is unstable.
  \item If $f'(x^*) = 0$, no conclusion is possible from $f'(x^*)$ alone.
\end{enumerate}
\end{theorem}

\begin{proof}
(a) By the definition of the derivative,
$\frac{f(x)}{x - x^*}\to f'(x^*)<0$, so there is $\delta>0$ with $f(x)<0$ on
$(x^*, x^*+\delta]$ and $f(x)>0$ on $[x^* - \delta, x^*)$. Take
$x(0)\in(x^*, x^*+\delta)$. The interval $J_0$ containing $x(0)$ has left
endpoint $x^*$, and the solution moves left, toward it. By
\cref{thm:ch07-structure}(b) it exists for all $t\ge0$, decreases, and tends
to $x^*$; in particular $|x(t) - x^*|\le|x(0) - x^*|$. The same holds from the
left. Given $\varepsilon>0$, take $\delta' = \min(\delta,\varepsilon)$: then
$|x(0) - x^*|<\delta'$ implies $|x(t) - x^*|<\varepsilon$ for all $t\ge0$, and
$x(t)\to x^*$. So $x^*$ is asymptotically stable.

(b) Choose $\delta>0$ with $f>0$ on $(x^*,x^*+\delta]$ (possible as in (a)).
For $x(0)\in(x^*,x^*+\delta)$, let $H(y) = \int_{x(0)}^{y}\frac{\mathrm{d}s}{f(s)}$.
Since $f\ge m>0$ on $[x(0),x^*+\delta]$, $H(x^*+\delta)\le\delta/m<\infty$; by
\cref{thm:ch03-separation}(c) the solution $x = H^{-1}(t)$ reaches $x^*+\delta$
at time $t = H(x^*+\delta)$. This holds however close $x(0)$ is to $x^*$, so
with $\varepsilon = \delta$ no $\delta'$ satisfies the definition of stability.

(c) $\dot x = -x^3$, $\dot x = x^3$, $\dot x = x^2$ all have $f(0) = f'(0) = 0$;
the origin is asymptotically stable, unstable, and \emph{semi-stable}
(attracting from one side, repelling from the other), respectively.
\end{proof}

% ------------------------------------------------------------
\section{Worked examples}\label{sec:ch07-examples}

\begin{example}[Reading a phase line]\label{ex:ch07-cubic}
Analyze $\dot x = f(x) = x(1-x)(x-2)$.

\emph{Solution.} Equilibria: $0,1,2$. Signs (product of three factors):
\begin{center}
\begin{tabular}{@{}lcccc@{}}
  \toprule
  interval & $x<0$ & $0<x<1$ & $1<x<2$ & $x>2$\\
  \midrule
  sign of $f$ & $+$ & $-$ & $+$ & $-$\\
  \bottomrule
\end{tabular}
\end{center}
So arrows point toward $0$ from both sides, away from $1$, and toward $2$.
Linearization agrees: $f'(x) = (1-x)(x-2) - x(x-2) + x(1-x)$ gives
$f'(0) = -2$, $f'(1) = 1$, $f'(2) = -2$. Hence $0$ and $2$ are asymptotically
stable and $1$ is unstable. By \cref{thm:ch07-structure}(b), a solution with
$x(0)\in(1,2)$ exists for all $t$, tends to $2$ as $t\to\infty$ and to $1$ as
$t\to-\infty$. The basin of attraction of $0$ is $(-\infty,1)$, and that of
$2$ is $(1,\infty)$ (\cref{fig:ch07-phase}).
\end{example}

\begin{figure}[ht]
  \centering
  % f(x) = x(1-x)(x-2) and the phase line below it (same x-scale).
\begin{tikzpicture}[>=Stealth]
\begin{axis}[deplot, name=graph, scale only axis, width=0.6\linewidth, height=0.32\linewidth, xmin=-0.6, xmax=2.6, ymin=-0.9, ymax=0.9,
  xlabel={$x$}, ylabel={$f(x)$}, ytick={-0.5,0.5}]
  \addplot[blue, samples=120, domain=-0.6:2.6] {x*(1-x)*(x-2)};
\end{axis}
\begin{axis}[at={(graph.south west)}, anchor=north west, yshift=-0.6cm,
  scale only axis, width=0.6\linewidth, height=1.2cm, xmin=-0.6, xmax=2.6, ymin=-1, ymax=1,
  axis line style={draw=none}, ticks=none, clip=false]
  \draw[thick] (axis cs:-0.6,0) -- (axis cs:2.6,0) node[right] {$x$};
  \draw[->, very thick, red!70!black] (axis cs:-0.5,0) -- (axis cs:-0.12,0);
  \draw[->, very thick, red!70!black] (axis cs:0.85,0) -- (axis cs:0.15,0);
  \draw[->, very thick, red!70!black] (axis cs:1.15,0) -- (axis cs:1.85,0);
  \draw[->, very thick, red!70!black] (axis cs:2.5,0) -- (axis cs:2.12,0);
  \fill (axis cs:0,0) circle (2.2pt);
  \node[below=3pt] at (axis cs:0,0) {$0$};
  \filldraw[fill=white] (axis cs:1,0) circle (2.2pt);
  \node[below=3pt] at (axis cs:1,0) {$1$};
  \fill (axis cs:2,0) circle (2.2pt);
  \node[below=3pt] at (axis cs:2,0) {$2$};
\end{axis}
\end{tikzpicture}

  \caption{Graph of $f(x) = x(1-x)(x-2)$ and the phase line of $\dot x = f(x)$
  below it. Filled dots: stable equilibria; open dot: unstable.}
  \label{fig:ch07-phase}
\end{figure}

\begin{example}[Needs care: the phase line does not show time]\label{ex:ch07-blowup}
The phase lines of $\dot x = x$ and $\dot x = x^2$ for $x>0$ look identical:
an arrow pointing to $+\infty$. But $x = x_0e^t$ exists for all $t$, while
$x = \frac{x_0}{1 - x_0t}$ blows up at $t = 1/x_0$. Whether $+\infty$ is reached
in finite time depends on whether $\int^\infty\frac{\mathrm{d}s}{f(s)}$
converges (\cref{prob:ch03-h3}), which an arrow cannot show.
\end{example}

\begin{example}[Needs care: a non-Lipschitz equilibrium is reached]\label{ex:ch07-reach}
For $\dot x = -x^{1/3}$ (real cube root), the phase line again shows arrows
toward $0$ from both sides. But $f$ is not $C^1$ at $0$, and
\cref{thm:ch07-structure}(a) fails: from $x(0) = 1$,
\notation{Leibniz, to separate}
$\int_x^1 s^{-1/3}\,\mathrm{d}s = t$ gives $\tfrac32\bigl(1 - x^{2/3}\bigr) = t$,
so $x$ reaches $0$ at $t = \frac32$ and can stay there (or, for the
time-reversed equation, leave it). ``Never reaches an equilibrium'' is a
theorem about $C^1$ (Lipschitz) right sides, not about phase lines.
\end{example}

\begin{example}[Harvesting: a saddle-node bifurcation]\label{ex:ch07-harvest}
Analyze \eqref{eq:ch07-harvest} completely.

\emph{Step 1: remove units.} \notation{Leibniz, for the change of variables.}
Put $X = x/K$ and $\tau = rt$. Then
$\der{X}{\tau} = \frac{1}{rK}\der{x}{t}$, and \eqref{eq:ch07-harvest} becomes
\begin{equation}\label{eq:ch07-harvest-dimless}
  \der{X}{\tau} = X(1 - X) - H,\qquad H = \frac{h}{rK}.
\end{equation}
Three parameters became one. (Always nondimensionalize before analysing: it
shows which combinations of parameters actually matter.)

\emph{Step 2: equilibria.} $X^2 - X + H = 0$ gives
$X_\pm = \frac{1\pm\sqrt{1 - 4H}}{2}$, real iff $H\le\frac14$.

\emph{Step 3: stability.} With $F(X) = X(1-X) - H$, $F'(X) = 1 - 2X$, so
$F'(X_+) = -\sqrt{1 - 4H}<0$ and $F'(X_-) = \sqrt{1-4H}>0$ for $H<\frac14$:
$X_+$ is asymptotically stable (the sustainable population), and $X_-$ is
unstable (a threshold: below it the stock declines). At $H = \frac14$ the two
merge into the semi-stable $X = \frac12$; for $H>\frac14$ there are no
equilibria and $F(X)\le\frac14 - H<0$ everywhere.

\emph{Step 4: collapse.} For $H>\frac14$, $\dot X\le -(H - \tfrac14)<0$, so by
the mean value theorem $X$ reaches $0$ in finite time from any start: the
fishery is extinct (the model stops being meaningful at $X = 0$).

\emph{Conclusion.} The maximum sustainable harvest is $h = rK/4$, at which the
population sits at $K/2$, a semi-stable state: any further loss sends it into
collapse. The change in the number of equilibria as $H$ crosses $\frac14$ is a
\emph{saddle-node bifurcation} (\cref{fig:ch07-bifurcation}); Chapter~22
studies bifurcations systematically.
\end{example}

\begin{figure}[ht]
  \centering
  % Equilibria of X' = X(1-X) - H: H = X(1-X).
\begin{tikzpicture}[>=Stealth]
\begin{axis}[deplot, axis lines=left, xmin=0, xmax=0.4, ymin=-0.1, ymax=1.1,
  xlabel={$H$}, ylabel={$X$}, height=0.45\linewidth,
  xtick={0,0.1,0.2,0.25,0.3,0.4}, xticklabels={$0$,$0.1$,$0.2$,$\frac14$,$0.3$,$0.4$}]
  \addplot[blue, very thick, samples=100, domain=0:0.25]
    {(1 + sqrt(1 - 4*x))/2};
  \addplot[blue, very thick, dashed, samples=100, domain=0:0.25]
    {(1 - sqrt(1 - 4*x))/2};
  \addplot[gray, dotted, samples=2] coordinates {(0.25,-0.1) (0.25,1.1)};
  % arrows at H = 0.1: up between branches, down outside
  \draw[->, red!70!black, thick] (axis cs:0.1,0.2) -- (axis cs:0.1,0.75);
  \draw[->, red!70!black, thick] (axis cs:0.1,1.05) -- (axis cs:0.1,0.93);
  \draw[->, red!70!black, thick] (axis cs:0.1,0.08) -- (axis cs:0.1,-0.07);
  % arrows at H = 0.33: always down
  \draw[->, red!70!black, thick] (axis cs:0.33,1.0) -- (axis cs:0.33,0.6);
  \draw[->, red!70!black, thick] (axis cs:0.33,0.5) -- (axis cs:0.33,0.1);
  \node[anchor=west, font=\small] at (axis cs:0.255,0.55) {collapse};
\end{axis}
\end{tikzpicture}

  \caption{Bifurcation diagram of $\dot X = X(1-X) - H$: equilibria
  $X_\pm(H)$. Solid: stable; dashed: unstable. Arrows show the direction of
  motion for fixed $H$. At $H = \frac14$ the branches meet and disappear.}
  \label{fig:ch07-bifurcation}
\end{figure}

% ------------------------------------------------------------
\section{Pitfalls}\label{sec:ch07-pitfalls}

\begin{pitfall}[Stability from the sign of $f$, not of $f'$]\label{pit:ch07-sign}
``$f(x^*) = 0$ and $f$ is decreasing near $x^*$'' means stable; ``$f(x^*)>0$''
means nothing, since $f(x^*) = 0$ at every equilibrium. Read the sign of $f$ on
\emph{either side} of $x^*$, or the sign of $f'(x^*)$.
\end{pitfall}

\begin{pitfall}[Treating $f'(x^*) = 0$ as stable]\label{pit:ch07-degenerate}
$f'(x^*) = 0$ is inconclusive: $\dot x = x^3$ is unstable and $\dot x = x^2$ is
semi-stable. Use the sign chart instead.
\end{pitfall}

\begin{pitfall}[Reading time off a phase line]\label{pit:ch07-time}
Arrows give direction, not speed. ``Tends to infinity'' may happen in finite
time (\cref{ex:ch07-blowup}), and ``tends to an equilibrium'' may happen in
finite time when $f$ is not Lipschitz (\cref{ex:ch07-reach}).
\end{pitfall}

\begin{pitfall}[Ignoring the physical domain]\label{pit:ch07-domain}
In \eqref{eq:ch07-harvest-dimless}, the region $X<0$ has arrows and solutions,
but no fish. A model's phase line is meaningful only on the set where the
model's assumptions hold.
\end{pitfall}

% ------------------------------------------------------------
\section{Problems}\label{sec:ch07-problems}

\subsection*{Warm-up}

\begin{problem}{W}\label{prob:ch07-w1}
Draw the phase line of $\dot x = x^2 - 4$ and classify the equilibria.
\end{problem}

\begin{problem}{W}\label{prob:ch07-w2}
Classify all equilibria of $\dot x = \sin x$.
\end{problem}

\begin{problem}{W}\label{prob:ch07-w3}
Classify the equilibria of $\dot x = x(1-x)^2$; for which one does
\cref{thm:ch07-linearization} fail to decide?
\end{problem}

\begin{problem}{W}\label{prob:ch07-w4}
Sketch the graph of a $C^1$ function $f$ for which $\dot x = f(x)$ has exactly
three equilibria: $-1$ semi-stable, $0$ unstable, $2$ asymptotically stable.
Is such an $f$ possible with all three equilibria satisfying
$f'(x^*)\neq0$?
\end{problem}

\begin{problem}{W}\label{prob:ch07-w5}
Show that the change of variables $X = x/K$, $\tau = rt$ turns
$\dot x = rx(1 - x/K) - Ex$ into an equation with a single parameter, and
identify it.
\end{problem}

\subsection*{Core}

\begin{problem}{C}\label{prob:ch07-c1}
Draw phase lines of $\dot x = rx - x^3$ for $r<0$, $r = 0$, $r>0$, and draw the
bifurcation diagram in the $(r,x)$ plane.
\end{problem}

\begin{problem}{C}\label{prob:ch07-c2}
Prove directly from \cref{thm:ch07-structure} that two distinct solutions of
$\dot x = f(x)$ ($f$ in $C^1$) that are both defined at time $t$ never take the
same value at that time, and that every solution is determined, up to a time
shift, by the value it takes at any one moment.
\end{problem}

\begin{problem}{C}\label{prob:ch07-c3}
A fishery uses \emph{proportional} harvesting:
$\dot x = rx(1 - x/K) - Ex$. Find the stable equilibrium for $0<E<r$, the
yield $Y(E) = Ex^*$, the effort $E$ maximizing it, and the maximum yield.
Compare the stability of the optimal state with \cref{ex:ch07-harvest}.
\end{problem}

\begin{problem}{C}\label{prob:ch07-c4}
(Allee effect.) For $\dot x = rx\bigl(\frac xA - 1\bigr)\bigl(1 - \frac xK\bigr)$
with $0<A<K$, draw the phase line, classify equilibria, and interpret $A$.
\end{problem}

\begin{problem}{C}\label{prob:ch07-c5}
For the logistic equation $\dot x = rx(1 - x/K)$, compute the time a solution
takes to go from $K/10$ to $9K/10$. Why does the answer not depend on the
initial condition?
\end{problem}

\begin{problem}{C}\label{prob:ch07-c6}
For \eqref{eq:ch07-harvest-dimless} with $H>\frac14$ and $X(0) = 1$, compute
exactly the time $\tau$ at which $X$ reaches $0$. What happens to this time as
$H\to\frac14^+$?
\end{problem}

\begin{problem}{C}\label{prob:ch07-c7}
For $\dot x = (x-1)^2 - \varepsilon$, draw the phase lines for $\varepsilon<0$,
$\varepsilon = 0$, $\varepsilon>0$. Explain why a semi-stable equilibrium is
``fragile'' under perturbation of $f$, while a hyperbolic one ($f'(x^*)\neq0$)
is not.
\end{problem}

\begin{problem}{C}\label{prob:ch07-c8}
\notation{Newton, for mechanics.} A body moves vertically with
$m\dot v = mg - kv|v|$ (down positive). Draw the phase line on the whole $v$-axis,
and prove that every solution tends to $v_T = \sqrt{mg/k}$.
\end{problem}

\subsection*{Hard}

\begin{problem}{H}\label{prob:ch07-h1}
Let $f$ be $C^1$, $f(x^*) = 0$, $f'(x^*)<0$. Prove there are $c,\delta>0$ such
that $|x(0) - x^*|<\delta$ implies
$|x(t) - x^*|\le|x(0) - x^*|\,e^{-ct}$ for all $t\ge0$.
\end{problem}

\begin{problem}{H}\label{prob:ch07-h2}
Let $f$ be merely \emph{continuous}. Solutions of $\dot x = f(x)$ need not be
unique. Prove that, nevertheless, every solution is monotone (not
necessarily strictly).
\end{problem}

\begin{problem}{H}\label{prob:ch07-h3}
(Hysteresis.) For $\dot x = h + x - \frac{x^3}{3}$, find the values of $h$ at
which the number of equilibria changes, draw the bifurcation diagram in the
$(h,x)$ plane, and describe what happens to a system sitting at the stable
equilibrium as $h$ is increased very slowly from $-2$ to $2$ and then decreased
back to $-2$.
\end{problem}

\begin{problem}{H}\label{prob:ch07-h4}
(Comparison principle.) Let $f$ be $C^1$, $x$ a solution of $\dot x = f(x)$ and
$u$ a differentiable function with $\dot u\le f(u)$ on $[0,T]$ and
$u(0)\le x(0)$. Prove $u(t)\le x(t)$ on $[0,T]$.
\end{problem}

\begin{problem}{H}\label{prob:ch07-h5}
(Ghost of a saddle-node.) Show that for $\varepsilon>0$, a solution of
$\dot x = \varepsilon + x^2$ takes time exactly $\pi/\sqrt\varepsilon$ to go from
$-\infty$ to $+\infty$. Then show that in \cref{prob:ch07-c6}, the collapse time
$\tau(H)$ satisfies $\tau(H)\sqrt{H - \tfrac14}\to\pi$ as $H\to\frac14^+$, and
explain the connection.
\end{problem}

\subsection*{Putnam/olympiad-style}

\begin{problem}{P}[Putnam 2009 B5, \cite{Putnam2009}]\label{prob:ch07-p1}
Let $f:(1,\infty)\to\R$ be a differentiable function such that
\[
  f'(x) = \frac{x^2 - f(x)^2}{x^2\bigl(f(x)^2 + 1\bigr)}\qquad\text{for all } x>1.
\]
Prove that $\displaystyle\lim_{x\to\infty}f(x) = \infty$.
\end{problem}

\begin{problem}{P}[Putnam 2010 B5, \cite{Putnam2010}]\label{prob:ch07-p2}
Is there a strictly increasing function $f:\R\to\R$ such that
$f'(x) = f(f(x))$ for all $x$?
\end{problem}

\begin{problem}{P}\label{prob:ch07-p3}
Let $g:\R\to\R$ be $C^1$ and periodic with period $1$, and let $x$ solve
$\dot x = g(x)$ for all $t\ge0$. Prove that $\lim_{t\to\infty}x(t)/t$ exists and
compute it in terms of $g$.
\end{problem}

% ------------------------------------------------------------
\section{Hints}\label{sec:ch07-hints}

\paragraph{Warm-up and core.}
\cref{prob:ch07-w4}: what does a sign chart with three zeros look like, and how
must the signs alternate if every zero is hyperbolic?
\cref{prob:ch07-c2}: if two solutions agreed at one time, use
\cref{thm:ch03-separation}(c).
\cref{prob:ch07-c5}: use the explicit solution of \cref{ex:ch03-logistic} and
\cref{prop:ch07-translation}.
\cref{prob:ch07-c6}: complete the square in $H - X + X^2$.
\cref{prob:ch07-c8}: $f(v) = g - \frac kmv|v|$ is $C^1$; find its only zero.

\hintfor{prob:ch07-h1}
\begin{hintlevels}
  \item Choose $\delta$ so that $\frac{f(x)}{x - x^*}\le -c$ on
        $0<|x - x^*|<\delta$ (with $c = \frac12|f'(x^*)|$).
  \item Differentiate $w(t) = (x(t) - x^*)^2$.
  \item $\dot w\le -2cw$; multiply by $e^{2ct}$ as in \cref{lem:ch02-exp}.
\end{hintlevels}

\hintfor{prob:ch07-h2}
\begin{hintlevels}
  \item Suppose $x(t_1)<x(t_2)$ and $x(t_2)>x(t_3)$ with $t_1<t_2<t_3$.
  \item Then $x$ takes some value $c$ in between both while increasing and while
        decreasing; look at $\dot x$ there.
  \item Find points where $\dot x = f(c)>0$ and $\dot x = f(c)<0$ for the same
        $c$; make this precise with the mean value theorem on suitable
        subintervals.
\end{hintlevels}

\hintfor{prob:ch07-h3}
\begin{hintlevels}
  \item The number of equilibria changes where $f = f' = 0$.
  \item $f' = 1 - x^2$; compute $h$ at $x = \pm1$.
  \item The stable branches are the outer ones; the state jumps when its branch
        disappears.
\end{hintlevels}

\hintfor{prob:ch07-h4}
\begin{hintlevels}
  \item Let $w = u - x$ and suppose $w(t_2)>0$ for some $t_2$.
  \item On an interval where $w>0$, $\dot w\le f(u) - f(x)\le Lw$ for a
        Lipschitz constant $L$ on the relevant bounded range.
  \item Take the last time $t_1<t_2$ with $w(t_1) = 0$ and show
        $\der{}{t}\bigl(e^{-Lt}w\bigr)\le0$ on $[t_1,t_2]$.
\end{hintlevels}

\hintfor{prob:ch07-h5}
\begin{hintlevels}
  \item \notation{Leibniz.} Time $= \int_{-\infty}^{\infty}\frac{\mathrm{d}x}{\varepsilon + x^2}$.
  \item In \eqref{eq:ch07-harvest-dimless}, put $Y = \tfrac12 - X$: then
        $\dot Y = (H - \tfrac14) + Y^2$.
  \item Most of the time is spent near $Y = 0$, where the equation looks like
        the first part; the distance travelled outside a small window costs
        bounded time.
\end{hintlevels}

\hintfor{prob:ch07-p1}
\begin{hintlevels}
  \item Compare $|f(x)|$ with $x$: where $|f(x)|<x$, $f'>0$.
  \item Show $f'(x)>-1/x^2$ for all $x$, so $f$ is bounded below on every
        $[x_0,\infty)$. If $f(x)<x$ somewhere, deduce that $|f(x)|<x$ holds at
        some point and then for all larger $x$ (the curves $f = \pm x$ cannot be
        crossed outward; compare slopes there).
  \item Then $f$ is increasing with a limit $L\in(-\infty,\infty]$; if $L$ were
        finite, what would $f'$ tend to?
\end{hintlevels}

\hintfor{prob:ch07-p2}
\begin{hintlevels}
  \item Explore what the equation forces: is $f'$ monotone? How does $f$
        behave as $x\to+\infty$?
  \item Show $f'$ is increasing and $f(x)\to\infty$; then that $f(x)>x$ for
        large $x$, so that $f'(x) = f(f(x))>f(x)$ eventually.
  \item Feed the resulting lower bound for $f$ back into $f'(x) = f(f(x))$ and
        compare with the blow-up criterion of \cref{prob:ch03-h3}.
\end{hintlevels}

\hintfor{prob:ch07-p3}
\begin{hintlevels}
  \item If $g$ has a zero, what are all solutions like?
  \item If $g>0$ (or $g<0$), the time to advance by one period is
        $T = \int_0^1\frac{\mathrm{d}s}{g(s)}$, independent of the start.
  \item Use \cref{prop:ch07-translation} and periodicity to show
        $x(t + T) = x(t) + 1$.
\end{hintlevels}

% ------------------------------------------------------------
\section{Further reading}\label{sec:ch07-reading}

\begin{itemize}
  \item \textcite[Ch.~2 and \S3.1]{Strogatz}: flows on the line, linear
        stability analysis, and the saddle-node bifurcation; the clearest
        introduction to this material.
  \item \textcite[\S\S1.2--1.3]{HSD}: the logistic model, constant harvesting,
        and bifurcations.
  \item \textcite[\S2.5]{BoyceDiPrima}: autonomous equations and population
        dynamics, including thresholds.
  \item \textcite[\S4.3]{Strogatz}: bottlenecks and the ``ghost'' of a
        saddle-node (\cref{prob:ch07-h5}).
  % VERIFY: Strogatz 2nd ed. §4.3 ("Nonuniform oscillator") discusses ghosts and bottlenecks.
\end{itemize}
